% !TEX root = Tesi_Corsi_Leonardo_656276.tex
\chapter{Conclusioni e sviluppi futuri}
\label{ch:conclusioni}

Il lavoro consegna due risultati, uno sull'array simulato e uno sul simulatore
che lo esegue. Il primo riguarda il canale ed è il principio enunciato nel
capitolo~\ref{ch:introduzione}; il secondo riguarda il ciclo di clock
parallelizzato e ne individua il regime di convenienza. Le due sezioni che
seguono li riprendono uno per uno, la terza indica gli sviluppi che le misure
lasciano aperti.

\section{Il \emph{ready bit} al posto della schedulazione}
\label{sec:sintesi-finale}

Il principio enunciato nel capitolo~\ref{ch:introduzione} è verificato nella
forma seguente. Sullo stesso programma e con gli stessi parametri, il risultato
ottenuto con il \emph{ready bit} è corretto per ogni velocità relativa dei nodi e
per ogni lunghezza della catena; senza, è corretto soltanto quando il consumatore
è abbastanza veloce, e negli altri casi l'array produce un risultato
numericamente definito ma errato, che nessun componente rileva. Ciascun lato del
canale commuta soltanto il proprio registro da un bit ed è il confronto fra i
due, e non il valore dell'uno o dell'altro, a stabilire se il canale sia pieno o
vuoto: nessuno dei due lati conosce il comportamento temporale dell'altro, e la
correttezza cessa di dipendere dalla schedulazione con cui l'host alimenta
l'array.

Il costo si misura su due assi. In cicli è nullo: da ritardo uno in poi le due
modalità impiegano lo stesso numero di cicli, perché il tempo lo
consuma il consumatore lento e non il protocollo; a ritardo zero il
\emph{lockstep} risparmia otto cicli su quarantacinque, ed è anche l'unica
configurazione in cui non perde dati, cioè l'unica in cui il \emph{ready bit}
non serviva. In area lo quantifica la sintesi del
modello RTL (sezione~\ref{sec:rtl}): 66 flip-flop e 6 porte
combinatorie per canale, di cui 64 flip-flop sono i due registri dati da 32 bit
e due sono i contatori, mentre le sei porte sono lo XOR e lo XNOR dei due
comparatori più quattro abilitazioni. Lo stato di sincronizzazione in senso
stretto è quindi di due flip-flop per canale, mentre il secondo registro dati è
una conseguenza dell'asimmetria fra \texttt{OUT} e \texttt{SETRDY}, poiché il
valore appena caricato e quello pubblicato e non ancora letto possono differire
nello stesso istante. Fra due celle adiacenti corrono 68 fili.

\section{L'efficienza della parallelizzazione dipende dalla sola dimensione
         della griglia}

Il secondo risultato riguarda il simulatore. Il ciclo di clock si parallelizza
senza una sola sezione critica, e il guadagno che se ne ottiene è determinato
dal solo prodotto $R \cdot C$: il lavoro di un ciclo sono $R \cdot C$ istruzioni
simulate, il costo da pagare per distribuirlo sono le due regioni parallele che
il ciclo apre, e il pareggio cade dove le due grandezze si eguagliano.

Il pareggio è misurato a $40\times40$, e la sezione~\ref{sec:pareggio} mostra
che la collocazione è stretta perché cade fra due forme contigue dello sweep
per tutti e tre i valori dell'intensità di calcolo. La previsione analitica lo
collocava a $24\times24$, ma lo scostamento non è del modello: la previsione
assume un costo per istruzione simulata misurato su una griglia molto più
grande, dove la località è peggiore, e rifatta con il costo della forma giusta
colloca il pareggio dove è stato misurato. Il modello analitico e la misura
concordano quindi sia sul rapporto a $12\times12$ sia sul punto di pareggio.

Che l'intensità di calcolo non compaia in quel rapporto discende dal disegno del
ciclo di clock, che è per costruzione \emph{una istruzione per cella}: il numero
di cicli si semplifica nel rapporto fra tempo seriale e tempo parallelo. Il
programma \texttt{pesante}
(sezione~\ref{sec:kernel-pesante}) verifica questa proprietà dove sarebbe più
visibile smentirla: trentadue volte le istruzioni di calcolo per ogni
comunicazione non spostano il pareggio di una forma, e il controllo a parità di
cicli non lascia un effetto residuo che la misura risolva. La misura del
pareggio si estende quindi ai programmi in cui le celle terminano insieme: per
questi, al di sopra di $40\times40$ l'aggiunta di thread accelera il ciclo di
clock qualunque cosa stiano eseguendo, e non occorre rimisurare il pareggio per
ognuno. Resta fuori il caso della terminazione sfalsata, perché la fase di
calcolo salta le celle già ferme e il lavoro per ciclo scende allora sotto
$R \cdot C$ distribuendosi male su una ripartizione statica; entrambi i
programmi su cui il pareggio è misurato hanno terminazione uniforme, quindi la
suite non lo esercita. La conseguenza
operativa è che per spostarlo non serve un programma diverso, serve una griglia
più grande oppure una regione parallela che non si riapra a ogni ciclo.

Le forme della suite di verifica si fermano a $12\times12$ e cadono quindi tutte
al di sotto del pareggio, dove la perdita arriva a poco più di un decimo della
velocità seriale. Sono le dimensioni che il compito di quello sweep richiede,
per la ragione data nella sezione~\ref{sec:metodologia}: quello sweep verifica
la correttezza e non le prestazioni.

Il tetto asintotico non è quello di Amdahl ma la località. La frazione
seriale scende dal $10{,}1\,\%$ su $8\times8$ allo $0{,}64\,\%$ su
$256\times256$, cioè si allontana esattamente nella direzione in cui l'array
diventa conveniente; nella stessa direzione il costo di simulare una istruzione
più che quintuplica, da 15,9 a 83,4\,ns, a un solo thread e a parità di lavoro
da svolgere, perché il passo di 16\,KB fra celle contigue disperde lo stato caldo
di ciascuna su una pagina diversa.

Il limite di queste misure è quello posto nella sezione~\ref{sec:pareggio}: la
macchina non è dedicata e il valore assoluto dello \emph{speedup} non è
riproducibile fra sessioni diverse. Si ritrovano invece in ogni sessione la
collocazione del pareggio, la severità della perdita al di sotto di esso,
l'andamento della frazione seriale e l'indipendenza dall'intensità di calcolo,
che sono le quattro affermazioni su cui il risultato poggia.

\section{Sviluppi futuri}
\label{sec:sviluppi}

Gli sviluppi possibili si dispongono su due piani. Le prime due modifiche non
sono scelte di progetto ma conseguenze dirette delle misure del
capitolo~\ref{ch:risultati}, e restano dentro il modello attuale. Le altre tre
lo estendono e ne mettono in discussione altrettante ipotesi, cioè la
profondità unitaria dei canali, il clock unico e il confine oltre il quale il
modello RTL si ferma.

\subsection{Le due modifiche indicate dalla misura}

Le prime due estensioni non sono scelte di progetto ma indicazioni che i dati
del capitolo~\ref{ch:risultati} lasciano aperte. Spostare il contatore di
lettura nella cella consumatrice ridurrebbe il \emph{false sharing} e
allineerebbe al tempo stesso il codice C alla struttura fisica del circuito,
dove quel flip-flop starebbe nel consumatore: due criteri indipendenti indicano
la stessa modifica, e non costa area, perché il contatore passa da un modulo
all'altro senza cambiare di numero. Non lo eliminerebbe, per la ragione data
nella sezione~\ref{sec:parallelo}, e richiede di decidere a chi spetti
l'aggiornamento di fine ciclo di quel campo, che oggi è del proprietario del
canale. Spostare la regione parallela fuori dal ciclo dei passi pagherebbe due
barriere invece di due per ciclo, ed è il modo di abbassare il pareggio senza
ingrandire la griglia.

\subsection{Canali di profondità maggiore di uno}

Portare la profondità da uno a $N$ è una estensione ortogonale: i due contatori
diventerebbero modulo $N$ invece che modulo due e il resto del protocollo non
cambierebbe di una riga. La misura da fare è quanta parte della differenza fra
le due modalità sopravvive quando il buffer nasconde il disallineamento.

\subsection{Clock indipendenti}

Il modello è sincrono, con un clock globale. Un modello asincrono vero
richiederebbe di cambiare il solo ciclo di simulazione e non il protocollo,
perché il \emph{ready bit} compensa il disallineamento del carico di lavoro e
non quello del clock.

\subsection{Il resto della cella in RTL}

Oggi il modello hardware copre il canale e l'interfaccia di comunicazione.
Estenderlo al processore completerebbe il confronto e permetterebbe di misurare
l'area dell'intera cella e non della sola interfaccia.
